\section{附录二：AI Lab 面对 Agent 范式的行动清单}

这期访谈可以转成一份面向模型团队的行动清单。它不是要求每家公司照抄，而是帮助读者把“范式变化”落到可执行问题上。

\noindent\begin{longtable}{p{0.25\textwidth}p{0.34\textwidth}p{0.32\textwidth}}
\toprule
决策项 & 需要回答的问题 & 错误做法 \\
\midrule
Agent 框架 & 我们的模型如何接入工具、环境、记忆和长程任务？ & 只把 Agent 当成复杂 system prompt。 \\
RL infra & rollout、reward、verifier、eval 如何规模化？ & 只做离线偏好数据，不建设环境闭环。 \\
卡资源 & research、pre-train、post-train 如何分配？ & 继续让预训练独占所有战略资源。 \\
模型家族 & 哪些环节用大模型，哪些环节用小模型或专用模型？ & 所有任务无差别调用最大模型。 \\
成本速度 & 端到端完成率、延迟和价格如何共同优化？ & 只看 benchmark，不算真实任务成本。 \\
组织结构 & 预训练、后训练、产品、系统是否共同进入闭环？ & 让单一团队垄断所有判断。 \\
价值观 & 替代哪些工作、保留哪些人类价值？ & 只优化破坏性或短期指标。 \\
\bottomrule
\end{longtable}

\begin{importantbox}{行动清单的用途}
如果一个团队声称要 all in Agent，但无法回答上表问题，它很可能只是在追热点。真正的 Agent 转型，会反映在资源分配、训练基础设施、模型路由、组织协作和价值观选择上。
\end{importantbox}

\subsection{本章小结}

Agent 范式不是一句战略口号，而是一组具体选择。读者可以用这张行动清单检查任何模型实验室、创业公司或产品团队是否真的进入了 Agent 时代。


